\begin{explanationbox}
	\textbf{\textcolor{CExplanation}{一句话直觉}}

	点差法通过设点作差，直接建立中点坐标与弦斜率的关系，避免联立解交点。

	\medskip
	\textbf{\textcolor{CExplanation}{核心拆解}}
	\begin{description}[style=nextline, leftmargin=2.8em, labelsep=0.5em, itemsep=0.45em]
		\item[\textbf{要点一（设点不求值）：}] 设弦两端点坐标为 \((x_1,y_1)\) 和 \((x_2,y_2)\)，但不需要求出具体数值。
		\item[\textbf{要点二（平方差消元）：}] 将两点坐标代入曲线方程并作差，利用平方差公式消去二次项，得到斜率与中点坐标的关系式。
	\end{description}

	\medskip
	\textbf{\textcolor{CExplanation}{几何本质（共轭性质）}}

	中点弦的斜率与中点坐标沿曲线的共轭方向，反映椭圆、双曲线、抛物线的中点与斜率的一一对应。

	\medskip
	\textbf{\textcolor{CExplanation}{代数意义（线性化）}}

	点差法将二次曲线的弦中点问题转化为一次方程的求解，使问题简便。

	\medskip
	\textbf{\textcolor{CExplanation}{考点价值（速算价值）}}
	\begin{itemize}[leftmargin=*, itemsep=0.5em, topsep=0.4em]
		\item \textbf{考法一（求中点弦方程）：} 已知弦的中点坐标，直接代入公式得斜率，再用点斜式写出直线方程。
		\item \textbf{考法二（求中点轨迹）：} 已知弦斜率满足条件，通过点差法得到中点坐标的方程，即中点轨迹。
	\end{itemize}

	\medskip
	\textbf{\textcolor{CExplanation}{顿悟点}}

	点差法不需求判别式，不涉及二次方程根的细节。

	\medskip
	\textbf{\textcolor{CExplanation}{使用场景}}
	\begin{itemize}[leftmargin=*, itemsep=0.5em, topsep=0.4em]
		\item \textbf{场景一（中点弦直接求）：} 题目直接给出弦的中点坐标，要求弦所在直线方程。
		\item \textbf{场景二（对称性问题）：} 涉及中点或对称关系时，利用点差法转化条件。
	\end{itemize}
\end{explanationbox}
